\section{Related Work}
\label{sec:related}
Table~\ref{tab:nearest} summarizes selected studies that connect agents or
software services to industrial equipment, production data, or process
operations. The studies locate the integration problem addressed here; they
are not matched performance baselines.

\subsection{Industrial Agents and Plant Integration}
Industrial-agent studies span process supervision, control generation, and
equipment-facing assistance. Large language model (LLM) frameworks address
end-to-end automation and industrial cyber-physical-system control
\cite{xia2025control,huang2026made}; related work generates programmable
logic controller (PLC) code or
combines human guidance with control
\cite{liu2026agents4plc,wu2025instructmpc}. Other studies examine resource
capabilities, shop-floor coordination, and digital-twin fault handling
\cite{lim2025llmenabled,zhao2026shopfloor,gill2025llmdt}, as well as
natural-language machine access and robotic-line control
\cite{hofmann2025nlc,waseem2026pretrained}. Biskupovic et al. are especially
close in using agents to supervise industrial systems: their agents select
services from a component library using performance and computational-load
metrics, monitor online performance, and reconfigure service selection when
performance degrades \cite{biskupovic2025supervision}. Taken together, this
strand studies agent behavior at the level of control, service choice, or
equipment interaction; the question here is how those agents connect to
scene-specific signals and governed operations.

Xia et al. connect cognitive, digital-twin, and physical-automation layers and
organize twin information as state snapshots, event messages, and plan
sequences \cite{xia2025architecture}. Other nearby designs expose factory-module
operations through the Model Context Protocol (MCP) with state-based tool
filtering, constrain manufacturing plans, or compose semantic recipes from
OPC Unified Architecture (OPC UA) skills
\cite{koehle2026llmdriven,silva2025mcpmanufacturing,koehle2025semanticrecipes}.
Several recent systems also combine agent execution with approval, industrial
data, or audit records. Graul et al. ground proposals in the Asset
Administration Shell (AAS) and route approved actions through OPC UA Methods,
with signed approvals and audit logging \cite{graul2026aas}. Schall separates
advisory language-model agents from process and instrumentation diagram
(P\&ID)-grounded verification and execution, and evaluates structural checks
across five simulated process scenarios \cite{schall2026safe}; its context
includes supervisory control and data acquisition (SCADA) and standard
operating procedure (SOP) information. Lin and Chen build a governed digital
thread across sensors, the manufacturing execution system (MES), and enterprise
resource planning (ERP) for closed-loop monitoring and control on real
manufacturing records \cite{lin2026agentic}. These systems address
complementary parts of industrial-agent deployment, including process
reasoning, service adaptation, asset grounding, action validation, and
evidence organization. This paper addresses a related but distinct integration
boundary: how to connect replaceable agent harnesses and task roles to
different configured production scenes while preserving
line--product--recipe--run context, semantic signal mappings, and a governed
parameter-write path. Whereas Biskupovic et al. study agent-based service
selection and reconfiguration, this architecture centers on scene configuration
and a shared operation boundary that resolves configured device nodes and
retains command receipts separately from process observations. The contribution
is an integration architecture and implementation, not a new process-control
law. Table~\ref{tab:nearest} compares reported mechanisms and evaluation
settings only; it does not rank features or performance.
\begin{table*}[!t]
\centering
\caption{Reported operation mechanisms and evaluations in selected
plant-facing studies. Evaluation settings and outcomes are source-specific;
this table does not compare performance across studies.}
\label{tab:nearest}
\small
\setlength{\tabcolsep}{3pt}
\begin{tabularx}{\textwidth}{@{}>{\raggedright\arraybackslash}p{1.8cm}>{\raggedright\arraybackslash}p{7.0cm}>{\raggedright\arraybackslash}X@{}}
\toprule
System & Reported operation or integration mechanism & Reported evaluation \\
\midrule
Xia et al. \cite{xia2025architecture} & Cognitive, twin, and physical layers; twin snapshots, event messages, and plan sequences & Robotic automation and process-simulation cases \\
Biskupovic et al. \cite{biskupovic2025supervision} & Agent-based selection and online reselection of services using performance and computational-load metrics & Reported adaptability and performance improvement in service-oriented industrial cyber-physical systems \\
K{\"o}hle et al. \cite{koehle2026llmdriven} & Injected live state and state-based filtering; MCP-to-OPC UA operations and Message Queuing Telemetry Transport (MQTT) coordination & Nine simulated production challenges \\
Vieira da Silva et al. \cite{silva2025mcpmanufacturing} & Constrained multi-step plans over MCP manufacturing functions & Laboratory multi-step process \\
K{\"o}hle et al. \cite{koehle2025semanticrecipes} & Semantic recipes composed from OPC UA skills; virtual commissioning & Unity assembly-workstation proof of concept \\
Graul et al. \cite{graul2026aas} & AAS-grounded proposals, signed human approval, PLC-validated OPC UA Methods, and immutable audit log & Two laboratory PoCs; single scripted runs \\
Lin and Chen \cite{lin2026agentic} & Governed sensor--MES--ERP digital thread with SOP-, risk-, and capability-aware action gating & 25,275 real records; operational metrics \\
Schall \cite{schall2026safe} & P\&ID-, SCADA-, and SOP-grounded proposal verification with gated execution & Five simulated scenarios; 43 proposals and 50 LLM runs \\
\midrule
This work & Common task-operation interface with registered harness adapters; scene-bound context and acquisition/write-node mappings; bounded named operations; separate command and observation records & One deterministic integrated simulator run; distinct protocol and guard checks, fixed-seed objective cases, and a scripted BOPET integration trace \\
\bottomrule
\end{tabularx}
\end{table*}

Industrial-informatics architectures offer complementary design precedents.
Gehrmann and Gunnarsson define trust boundaries and synchronization
requirements for digital-twin-based industrial control and evaluate a
protected-update proof of concept
\cite{gehrmann2020dtiacs}. Mahmood et al. organize special-purpose industrial
IoT functions around network requirements, while Sengupta et al. distribute
security services across fog and device layers and report simulation and
testbed evidence \cite{mahmood2023functional,sengupta2021securefog}. These
works address communication, security, and functional decomposition. The
present system focuses on a different interface: binding scene-specific
context and device mappings to task-level operations while retaining
execution-time checks and separate command and process records. Its operation
checks complement safety-aware decision support
\cite{jiang2026agentic} and do not replace configured plant protection.

\subsection{Production Context, Agent Interfaces, and Automation}
Industrial integration also depends on abstractions that are not specific to
agent architectures. The Model Context Protocol (MCP) specifies interoperable
tool discovery and call mechanisms \cite{mcp2024}. The platform dispatcher's
operation-level scope, policy, grant, and binding checks are distinct from MCP
transport authorization. In parallel,
ISA-88 and ISA-95 describe recipes and production operations
\cite{isa88_2010,isa95_2025}. Run-to-run methods connect process observations
to successive production runs \cite{moyne2000r2r}. Human--automation research
also formalizes the allocation of information acquisition, analysis, decisions,
and actions between people and automation \cite{parasuraman2000model}. At the
agent level, AgentSpec and CaMeL address tool behavior and capability-based
information flow \cite{wang2025agentspec,debenedetti2025camel}. These
standards and concepts inform the design's treatment of production context,
tool interaction, human allocation, and information-flow boundaries; their
mention does not imply implementation of every listed standard or safety
mechanism. Within this selected literature, these concerns are treated through
separate interfaces and validation layers; how to bind scene context and
heterogeneous I/O to a shared task-operation interface and write-governance
path remains under-specified. This paper addresses that integration question,
and the evaluation tests selected links in a deterministic simulator run.
